\begin{lemma}
\label{lemma-complete-devissage-flat-finitely-presented-module}
Let $R \to S$ be a ring map of finite presentation.
Let $N$ be a finitely presented $S$-module flat over $R$.
Let $\mathfrak r \subset R$ be a prime ideal.
Assume there exists a complete d\'evissage of $N/S/R$ over $\mathfrak r$.
Then there exists an $f \in R$, $f \not \in \mathfrak r$
such that
$$
N_f \cong B_1^{\oplus r_1} \oplus \ldots \oplus B_n^{\oplus r_n}
$$
as $R$-modules where each $B_i$ is a smooth $R_f$-algebra with geometrically
irreducible fibres. Moreover, $N_f$ is projective as an $R_f$-module.
\end{lemma}

\begin{proof}
Let $(A_i, B_i, M_i, \alpha_i)_{i = 1, \ldots, n}$ be the given
complete d\'evissage. We prove the lemma by induction on $n$.
Note that the assertions of the lemma are entirely about the structure
of $N$ as an $R$-module. Hence we may replace $N$ by $M_1$, and we
may think of $M_1$ as a $B_1$-module. See
Remark \ref{remark-finite-presentation}
in order to see why $M_1$ is of finite presentation as a $B_1$-module. By
Lemma \ref{lemma-induction-step-fp}
we may, after replacing $R$ by $R_f$ for some
$f \in R$, $f \not \in \mathfrak r$, assume
the map $\alpha_1 : B_1^{\oplus r_1} \to M_1$ is $R$-universally injective.
Since $M_1$ and $B_1^{\oplus r_1}$ are $R$-flat and finitely presented as
$B_1$-modules we see that $\Coker(\alpha_1)$ is $R$-flat
(Algebra, Lemma \ref{algebra-lemma-ui-flat-domain})
and finitely presented as a $B_1$-module. Note that
$(A_i, B_i, M_i, \alpha_i)_{i = 2, \ldots, n}$ is a complete
d\'evissage of $\Coker(\alpha_1)$. Hence the induction hypothesis
implies that, after replacing
$R$ by $R_f$ for some $f \in R$, $f \not \in \mathfrak r$,
we may assume that $\Coker(\alpha_1)$ has a decomposition
as in the lemma and is projective. In particular
$M_1 = B_1^{\oplus r_1} \oplus \Coker(\alpha_1)$.
This proves the statement regarding the decomposition.
The statement on projectivity follows as $B_1$ is projective as
an $R$-module by
Lemma \ref{lemma-fibres-irreducible-flat-projective-nonnoetherian}.
\end{proof}
